\chapter{Clonezilla}

\section{Instalacion de Clonezilla en un pendrive}

En \href{Clonezilla Download}{https://clonezilla.org/downloads/download.php} buscar la última versión estable y seguis la guía para bajarse la imagen en .ISO

Luego existen varios metodos para pasar la imagen a un pendrive. Desde windows la que generalmente uso es \href{Rufus}{https://rufus.ie/es/}, solo selecciono el pendrive como destino, la imagen como origen y listo, le doy a START 

Luego ya se puede poner la PC para que arranque desde un usb como (entrando al BIOS), poner el pendrive con clonzilla y arrancar la PC.

\section{Creacion de Imagenes}
Lo primero para la creación de la imagen es tener una instalación completa (todo el software y drivers instalados y funcionando), límpia (sin uso por los alumnos) y debidamente probada (con scripts que hayan simulado el uso habitual, sobre todo de las instalaciones de Python).

Luego se pone en la PC el medio para guardar la imagen (puede ser un medio en blanco para no arriesgarse con la sobreescritura/borrado del dispositivo con todas las imagenes) y se inicia la PC con Clonzilla.

En el menu\dots
\todo[inline]{terminar la arte importante de la explicación}

\subsection{Imagen a Disco}
\todo[inline]{Creacion de imagenes}
